% Options for packages loaded elsewhere
\PassOptionsToPackage{unicode}{hyperref}
\PassOptionsToPackage{hyphens}{url}
\documentclass[
  ignorenonframetext,
]{beamer}
\newif\ifbibliography
\usepackage{pgfpages}
\setbeamertemplate{caption}[numbered]
\setbeamertemplate{caption label separator}{: }
\setbeamercolor{caption name}{fg=normal text.fg}
\beamertemplatenavigationsymbolsempty
% remove section numbering
\setbeamertemplate{part page}{
  \centering
  \begin{beamercolorbox}[sep=16pt,center]{part title}
    \usebeamerfont{part title}\insertpart\par
  \end{beamercolorbox}
}
\setbeamertemplate{section page}{
  \centering
  \begin{beamercolorbox}[sep=12pt,center]{section title}
    \usebeamerfont{section title}\insertsection\par
  \end{beamercolorbox}
}
\setbeamertemplate{subsection page}{
  \centering
  \begin{beamercolorbox}[sep=8pt,center]{subsection title}
    \usebeamerfont{subsection title}\insertsubsection\par
  \end{beamercolorbox}
}
% Prevent slide breaks in the middle of a paragraph
\widowpenalties 1 10000
\raggedbottom
\AtBeginPart{
  \frame{\partpage}
}
\AtBeginSection{
  \ifbibliography
  \else
    \frame{\sectionpage}
  \fi
}
\AtBeginSubsection{
  \frame{\subsectionpage}
}
\usepackage{iftex}
\ifPDFTeX
  \usepackage[T1]{fontenc}
  \usepackage[utf8]{inputenc}
  \usepackage{textcomp} % provide euro and other symbols
\else % if luatex or xetex
  \usepackage{unicode-math} % this also loads fontspec
  \defaultfontfeatures{Scale=MatchLowercase}
  \defaultfontfeatures[\rmfamily]{Ligatures=TeX,Scale=1}
\fi
\usepackage{lmodern}
\ifPDFTeX\else
  % xetex/luatex font selection
\fi
% Use upquote if available, for straight quotes in verbatim environments
\IfFileExists{upquote.sty}{\usepackage{upquote}}{}
\IfFileExists{microtype.sty}{% use microtype if available
  \usepackage[]{microtype}
  \UseMicrotypeSet[protrusion]{basicmath} % disable protrusion for tt fonts
}{}
\makeatletter
\@ifundefined{KOMAClassName}{% if non-KOMA class
  \IfFileExists{parskip.sty}{%
    \usepackage{parskip}
  }{% else
    \setlength{\parindent}{0pt}
    \setlength{\parskip}{6pt plus 2pt minus 1pt}}
}{% if KOMA class
  \KOMAoptions{parskip=half}}
\makeatother
\usepackage{longtable,booktabs,array}
\newcounter{none} % for unnumbered tables
\usepackage{calc} % for calculating minipage widths
\usepackage{caption}
% Make caption package work with longtable
\makeatletter
\def\fnum@table{\tablename~\thetable}
\makeatother
\setlength{\emergencystretch}{3em} % prevent overfull lines
\providecommand{\tightlist}{%
  \setlength{\itemsep}{0pt}\setlength{\parskip}{0pt}}
% Slide layout defaults for warning-clean scientific decks.
\usepackage{etoolbox}
\IfFileExists{xurl.sty}{\usepackage{xurl}}{}
\IfFileExists{seqsplit.sty}{\usepackage{seqsplit}}{\newcommand{\seqsplit}[1]{#1}}
\protected\def\breakseq#1{\seqsplit{#1}}
\protected\def\breaktt#1{\begingroup\ttfamily\seqsplit{#1}\endgroup}
\setlength{\emergencystretch}{6em}
\tolerance=5000
\hbadness=10000
\hfuzz=1pt
\setlength{\tabcolsep}{2pt}
\AtBeginEnvironment{longtable}{\tiny\renewcommand{\arraystretch}{0.86}\setlength{\tabcolsep}{1pt}}
\AtBeginEnvironment{tabular}{\tiny\renewcommand{\arraystretch}{0.86}\setlength{\tabcolsep}{1pt}}

% Natbib and cross-reference fallbacks — slides don't load natbib
% or cleveref, but combined-PDF manuscript prose may emit these
% commands. The fallback renders citations as a bracketed key list
% and cross-references as detokenized labels so slides don't fail on
% undefined control sequences or raw underscores. \providecommand is
% a no-op if packages load later.
\providecommand{\citep}[1]{[#1]}
\providecommand{\citet}[1]{#1}
\providecommand{\citealp}[1]{#1}
\providecommand{\citeauthor}[1]{#1}
\providecommand{\citeyear}[1]{#1}
\providecommand{\cref}[1]{\texttt{\detokenize{#1}}}
\providecommand{\Cref}[1]{\texttt{\detokenize{#1}}}

\newtheorem{proposition}{Proposition}
\newtheorem{hypothesis}{Hypothesis}
\usepackage{bookmark}
\IfFileExists{xurl.sty}{\usepackage{xurl}}{} % add URL line breaks if available
\urlstyle{same}
% fallback for those not using the hyperref driver hyperxmp:
\makeatletter
\@ifundefined{xmpquote}{\newcommand{\xmpquote}[1]{#1}}{}
\makeatother
\hypersetup{
  hidelinks,
  pdfcreator={LaTeX via pandoc}}

\author{\texorpdfstring{}{}}
\date{}

\begin{document}

\section{Appendix B: Technical Notes}\label{appendix-b-technical-notes}

\begin{frame}[fragile,allowframebreaks]{Determinism}
\protect\phantomsection\label{determinism}
All stochastic steps use fixed seeds (seed = 42 for NMF, SVD, and graph
layouts). The fixture corpus, TF-IDF/SVD embeddings, and topic
factorization are byte-stable across runs. Graph centrality scores are
rounded to a fixed precision and ranked with a node-id tiebreaker so
that floating-point non-associativity in iterative solvers cannot
perturb the reported rankings. Record identity uses a content digest
(SHA-1, \breaktt{usedforsecurity=False}) rather than a salted hash, so
de-duplication and corpus byte-stability hold across processes.
\end{frame}

\begin{frame}[fragile,allowframebreaks]{Data Model}
\protect\phantomsection\label{data-model}
Each record is a \texttt{Paper} dataclass with: title, abstract, authors
(list of \texttt{Author}), year, DOI, arXiv ID, Semantic Scholar ID,
OpenAlex ID, venue, citation count, references (list of canonical IDs),
publication date, PDF URL, open-access flag, and full-text source. The
canonical identifier hierarchy governs de-duplication and citation
resolution:

\[
\text{canonical\_id} = \begin{cases}
\texttt{doi:} + \text{normalize}(\text{DOI}) & \text{if DOI present} \\
\texttt{arxiv:} + \text{arXiv\_id} & \text{if arXiv ID present} \\
\texttt{s2:} + \text{S2\_id} & \text{if S2 ID present} \\
\texttt{openalex:} + \text{OpenAlex\_id} & \text{if OpenAlex ID present} \\
\texttt{title:} + \text{SHA1}(\text{title})[:16] & \text{otherwise}
\end{cases}
\]

DOI normalization lower-cases the DOI and strips any
\breaktt{https://doi.org/} or \texttt{dx.doi.org/} prefix, so the same
paper returned by two engines under different case or format variants
merges to a single canonical ID.
\end{frame}

\begin{frame}[allowframebreaks]{NMF Mathematics}
\protect\phantomsection\label{nmf-mathematics}
Non-negative matrix factorization decomposes the TF-IDF document-term
matrix \(\mathbf{V} \in \mathbb{R}^{m \times n}\) (where \(m\) is the
number of documents and \(n\) is the vocabulary size) into
\(\mathbf{W} \in \mathbb{R}^{m \times k}\) and
\(\mathbf{H} \in \mathbb{R}^{k \times n}\), where \(k\) is the number of
topics (here 5). The factorization minimizes:

\[
\min_{\mathbf{W}, \mathbf{H} \geq 0} \|\mathbf{V} - \mathbf{W} \mathbf{H}\|_F^2
\]

using multiplicative update rules \breakseq{[@lee1999learning]} with a fixed
random seed for reproducibility. The topic-term matrix \(\mathbf{H}\)
gives the top terms per topic; the document-topic matrix \(\mathbf{W}\)
gives each document's topic loadings.
\end{frame}

\begin{frame}[allowframebreaks]{Growth Rate Estimation}
\protect\phantomsection\label{growth-rate-estimation}
The compound annual growth rate is:

\[
\text{CAGR} = \left(\frac{N_{\text{end}}}{N_{\text{start}}}\right)^{1/(T_{\text{end}} - T_{\text{start}})} - 1
\]

where \(N_{\text{start}}\) and \(N_{\text{end}}\) are the publication
counts in the first and last years of the corpus, respectively. The
doubling time is \(t_d = \ln(2) / \ln(1 + \text{CAGR})\). For this run:
CAGR = 3.45\%, doubling time = 11.3 years.
\end{frame}

\begin{frame}[fragile,allowframebreaks]{Configuration Surface}
\protect\phantomsection\label{configuration-surface}
A single \breaktt{manuscript/config.yaml} controls the search term,
per-engine query and keyword sets, engine enable toggles, subfield
taxonomy, hypotheses, full-text and embedding options, and paper
metadata. This run drew on 7 engines, a 6-bucket taxonomy, and 6
hypotheses.
\end{frame}

\begin{frame}[fragile,allowframebreaks]{Artifacts}
\protect\phantomsection\label{artifacts}
Intermediate and final outputs live under \texttt{output/} and are
disposable and regenerable:

{\def\LTcaptype{none} % do not increment counter
\begin{longtable}[]{@{}
  >{\raggedright\arraybackslash}p{(\linewidth - 4\tabcolsep) * \real{0.3333}}
  >{\raggedright\arraybackslash}p{(\linewidth - 4\tabcolsep) * \real{0.3333}}
  >{\raggedright\arraybackslash}p{(\linewidth - 4\tabcolsep) * \real{0.3333}}@{}}
\toprule\noalign{}
\begin{minipage}[b]{\linewidth}\raggedright
File
\end{minipage} & \begin{minipage}[b]{\linewidth}\raggedright
Stage
\end{minipage} & \begin{minipage}[b]{\linewidth}\raggedright
Description
\end{minipage} \\
\midrule\noalign{}
\endhead
\texttt{corpus.jsonl} & 01 & De-duplicated corpus (2302 records) \\
\breaktt{temporal\_analysis.json} & 02 & Year counts, CAGR, doubling
time, peak year \\
\breaktt{subfield\_classification.json} & 02 & Per-bucket paper counts \\
\breaktt{subfield\_timeline.json} & 02 & Per-subfield annual breakdown \\
\texttt{tfidf\_data.json} & 02 & TF-IDF matrix, feature names, doc
tokens \\
\texttt{topics.json} & 02 & NMF topic-term distributions \\
\breaktt{citation\_network.json} & 02 & Network metrics, PageRank, HITS,
communities \\
\breaktt{citation\_graph.gml} & 02 & GraphML citation graph \\
\breaktt{nanopublications.jsonl} & 03 & LLM-extracted assertions (0 in
this run) \\
\breaktt{hypothesis\_scores.json} & 03 & Per-hypothesis evidence
scores \\
\breaktt{fulltext\_assessment.json} & 06 & Abstract/OA/PDF coverage
report \\
\bottomrule\noalign{}
\end{longtable}
}
\end{frame}

\end{document}
